\chapter{Limitations and Threats to Validity}
\label{ch:limitations}

A monograph that claims an ``immutable system of safety boundaries'' must be
honest about where the boundaries are not yet closed. This chapter enumerates
the limitations.

\section{Scaffolding vs. Implementation}
\label{sec:lim-scaffold}

Several layers are real but partial:
\begin{itemize}
  \item The orchestration glue in \texttt{ci/run\_e2e.sh} is commented;
        the layers it would connect exist individually but are not yet wired
        into a single command.
  \item The Liquid refinement predicates are specified (\texttt{todo} for
        non-literal cases) rather than enforced.
  \item \texttt{Core.lean} is a stub; the theorem library is the active
        frontier tracked by \texttt{sorrydb}.
\end{itemize}

\section{Certificate Stubs}
\label{sec:lim-cert}

In \texttt{verified\_kernel.cm}, the WORM-ledger validation inside
\texttt{rewrite\_with\_cert} is a stub. The \emph{interface} (a certificate
pointer threaded through every rewrite) is present, but the cryptographic
check it implies is not yet enforced in the running kernel. The security model
of \cref{ch:security} is therefore partly \emph{architectural} and not yet
\emph{operational} at the kernel's innermost loop.

\section{Equivalence, Not Proof}
\label{sec:lim-equiv}

The FFI cross-validation shows the C{--} kernel and Lean agree on tested
inputs. It does not prove the closed forms; proving them is Lean's
responsibility at Layer~4, which is currently placeholder for most goals. The
system's strongest guarantee today is \emph{consistency across independent
implementations}, not \emph{formal proof of every result}.

\section{Scope of the Frontend}
\label{sec:lim-apl}

The \textsc{apl} frontend parses a tractable fragment. Tacit \textsc{apl},
operator trains, and the full glyph set are out of scope. Inputs outside the
fragment are rejected, which preserves safety, but limits expressiveness.

\section{Threats to Validity}
\label{sec:lim-threats}

\begin{itemize}
  \item \textbf{Internal validity}: the benchmark numbers are
        hardware-dependent and not yet statistically characterized.
  \item \textbf{External validity}: the three closed-form patterns are a small
        sample; generalization to arbitrary symbolic mathematics is unproven.
  \item \textbf{Construct validity}: ``coherence'' and ``completeness'' are
        operationalized by simple metrics (entropy, sorry-ratio) that may not
        capture all relevant safety properties.
\end{itemize}

\section{Mitigations}
\label{sec:lim-mit}

Each limitation is recorded as a gap, not hidden:
\begin{itemize}
  \item Open goals live in \texttt{sorrydb}; the completeness score $C$
        makes progress visible.
  \item \texttt{metadata.json} status fields (\texttt{awaiting\_worm\_seal})
        mark unsealed artifacts.
  \item The build gate fails loudly on any broken boundary.
\end{itemize}
The discipline of inspectability turns limitations into tracked work items
rather than silent risks.
